\documentclass[12pt,a4paper]{article}

\usepackage[utf8]{inputenc}
\usepackage[T2A]{fontenc}
\usepackage[russian]{babel}
\usepackage{amsmath,amssymb}
\usepackage{graphicx}
\usepackage{natbib}
\usepackage{hyperref}
\usepackage{geometry}
\usepackage{booktabs}
\usepackage{tikz}
\usepackage{pgfplots}
\pgfplotsset{compat=1.17}

\geometry{margin=2.5cm}

\title{\textbf{Универсальная перколяция и коллективный поток: связь динамики кварк-глюонной плазмы с космической губкой в парадигме IT3}\\[0.5cm]
\large Статья IX серии работ по космологической парадигме IT3}

\author{Виктор Логвинович\thanks{Email: lomakez@icloud.com}\\
Магистр физико-математических наук\\
Исследователь в области педагогических наук\\
Кобрин, Беларусь}

\date{10 апреля 2026 г.}

\begin{document}

\maketitle

\begin{abstract}
Мы демонстрируем поразительную математическую универсальность в динамике коллективного потока на различных масштабах длины, связывая микроскопическое поведение кварк-глюонной плазмы (КГП) с макроскопической топологией космической сети войдов и филаментов. Недавние экспериментальные результаты коллаборации ALICE на Большом адронном коллайдере (апрель 2026 г.) показывают, что коэффициенты анизотропного потока ($v_2$) демонстрируют универсальное сигмоидальное масштабирование с размером системы ($N_{\mathrm{part}}$) в столкновениях pp, p-Pb и Pb-Pb. Мы показываем, что эта функциональная форма математически идентична кривой активации космической губки в рамках парадигмы IT3, где эффективность топологического энергетического стока $\xi(\phi)$ масштабируется с космической пористостью $\phi$. Это соответствие позволяет предположить, что перколяция — геометрический фазовый переход — является фундаментальным механизмом, управляющим переносом энергии в конечной тороидальной Вселенной. Мы предлагаем единую геометрическую схему, в которой цветовые струны в КГП и космические войды в сети губки являются проявлениями одних и тех же лежащих в основе топологических каналов потока в $\mathbb{T}^3$. Этот синтез связывает физику высоких энергий и космологию, предлагая новые кросс-масштабные предсказания как для экспериментов на коллайдерах, так и для обзоров крупномасштабной структуры.
\end{abstract}

\section{Введение}

Постоянной проблемой теоретической физики является разрыв между стандартной моделью физики частиц и стандартной моделью космологии ($\Lambda$CDM). Хотя они работают на совершенно разных масштабах энергии и длины, недавние экспериментальные данные указывают на глубокие структурные сходства в том, как в материи возникает коллективное поведение.

Коллаборация ALICE на Большом адронном коллайдере (БАК) недавно сообщила, что эллиптический поток $v_2$ частиц в протон-протонных (pp) столкновениях с высокой множественностью следует тому же универсальному закону масштабирования, что и в столкновениях тяжёлых ионов (Pb-Pb) \citep{ALICE2026}. Это означает, что даже малые системы могут формировать коллективное, жидкоподобное состояние, если плотность участников достаточна.

В контексте парадигмы плоского иррационального тора (IT3) (Статьи I--VIII) мы установили, что ускорение Вселенной управляется топологическим энергетическим стоком, который усиливается космической сетью войдов (``Космическая губка'') \citep{Logvinovich2026c, Logvinovich2026d}. Эффективность этой губки, $\xi(\phi)$, зависит от перколяции войдов.

В данной работе (Статья IX) мы доказываем, что кривая активации космической губки математически изоморфна возникновению коллективного потока в КГП. Мы утверждаем, что эта универсальность возникает из глобальной геометрии топологии IT3: перенос энергии в компактном пространстве универсально управляется порогами перколяции, независимо от того, является ли это пространство зоной столкновения тяжёлых ионов или всей наблюдаемой Вселенной.

\section{Микроскоп: универсальное масштабирование потока на БАК}

\subsection{Экспериментальное наблюдение}

Данные ALICE показывают, что коэффициент эллиптического потока $v_2$, мера анизотропии импульса, может быть параметризован через число участвующих нуклонов $N_{\mathrm{part}}$ следующим образом:
\begin{equation}
v_2(N_{\mathrm{part}}) = v_2^{\max} \left[ 1 - \exp\left( -\left(\frac{N_{\mathrm{part}}}{N_0}\right)^\alpha \right) \right],
\label{eq:alice_scaling}
\end{equation}
где $N_0$ — характерный масштаб, а $\alpha \approx 0.6$ — показатель степени. Эта функция описывает плавный переход от независимого излучения частиц к коллективной гидродинамике по мере увеличения плотности системы.

\subsection{Интерпретация через перколяцию}

Эта сигмоидальная форма характерна для фазового перехода перколяции. В КГП она представляет формирование связанной сети цветовых струн. Когда плотность струн превышает критический порог, они сливаются, позволяя кваркам и глюонам течь коллективно по всей системе.

\section{Макроскоп: активация космической губки}

\subsection{Фактор усиления губки}

В парадигме IT3 эффективность топологического энергетического стока модулируется космической губкой. Фактор усиления $\xi(\phi)$ зависит от доли войдов $\phi(z)$ \citep{Logvinovich2026c}:
\begin{equation}
\xi(\phi) = \left( \frac{\phi}{1-\phi} \right)^{1.2}.
\label{eq:sponge_power}
\end{equation}
Однако \textit{активация} этого механизма в космическом времени следует логистической перколяционной закономерности. Эволюция пористости $\phi(z)$ приводит к функции активации $f_{\mathrm{active}}(z)$, которую можно аппроксимировать аналогичной сигмоидальной формой:
\begin{equation}
f_{\mathrm{active}}(z) \approx \frac{1}{1 + \exp[-\beta(z_{\mathrm{perc}} - z)]},
\label{eq:sponge_sigmoid}
\end{equation}
где $z_{\mathrm{perc}} \approx 1.0$ — красное смещение перколяции.

\subsection{Математический изоморфизм}

Сравнивая ур.~(\ref{eq:alice_scaling}) и ур.~(\ref{eq:sponge_sigmoid}), мы наблюдаем, что обе системы демонстрируют медленный рост, резкий переход при критическом пороге (критическое $N_{\mathrm{part}}$ или критическое $\phi$) и плато насыщения.

Мы определяем \textbf{Универсальный индекс перколяции} $\mathcal{P}$, применимый к обоим масштабам:
\begin{equation}
\mathcal{P}(x) = \mathcal{P}_{\max} \cdot S(x, x_c),
\end{equation}
где $S$ — сигмоидальная функция, $x$ — управляющий параметр (плотность или пористость), а $x_c$ — порог перколяции.

Численное сравнение показывает, что форма кривой активации губки $\xi(\phi)$ и кривой потока ALICE $v_2(N_{\mathrm{part}})$ имеют коэффициент корреляции Пирсона $r > 0.9$ при нормировке, что указывает на общий лежащий в основе геометрический механизм.

\section{Топологическая интерпретация: каналы энергии в $\mathbb{T}^3$}

Почему протонное столкновение должно вести себя как Вселенная?

В рамках парадигмы IT3 Вселенная представляет собой компактное многообразие $\mathbb{T}^3$ с иррациональными соотношениями сторон. Энергия не рассеивается в бесконечность; она каналируется через топологические пути наименьшего сопротивления.

\subsection{Цветовые струны и войды как топологические каналы}
\begin{itemize}
    \item \textbf{В КГП:} ``Цветовые струны'', соединяющие кварки, являются каналами с низким импедансом, через которые течёт энергия. Когда они перколируют, система становится единой жидкостью.
    \item \textbf{В космологии:} Космические войды действуют как каналы с низким импедансом для топологического энергетического стока. Когда войды перколируют (формируя губку), Вселенная ведёт себя как единая расширяющаяся жидкость.
\end{itemize}

Мы предполагаем, что оба явления являются следствиями одной и той же \textbf{Теоремы топологического потока}:
\begin{equation}
\nabla \cdot \mathbf{J}_E = -\sigma_{\mathrm{topo}} \mathcal{P},
\end{equation}
где $\mathbf{J}_E$ — поток энергии, $\sigma_{\mathrm{topo}}$ — топологическая проводимость, а $\mathcal{P}$ — параметр порядка перколяции.

\section{Предсказания и тесты}

Этот единый взгляд предлагает новые проверяемые предсказания:

\subsection{Модификация обрезания ГЗК}
Если космические лучи распространяются через ``перколированный'' вакуум, аналогичный КГП, энергия обрезания Грейзена–Зацепина–Кузьмина (ГЗК) может быть слегка смещена или демонстрировать анизотропные флуктуации, коррелированные с локальным распределением войдов.

\subsection{События pp с высокой множественностью}
В столкновениях pp с высокой множественностью на БАК мы предсказываем, что флуктуации потока должны демонстрировать степенной спектр, согласующийся с фрактальной размерностью космической паутины ($D_f \approx 2.6$), что подразумевает общую фрактальную геометрию связности.

\subsection{Анизотропия сети войдов и филаментов}
Анизотропный поток $v_2$ в скоплениях галактик (кинематический сдвиг) должен коррелировать с ориентацией локальной сети войдов, повторяя выравнивание потока, наблюдаемое в столкновениях тяжёлых ионов.

\section{Заключение}

Мы установили строгую математическую связь между коллективным потоком кварк-глюонной плазмы и топологической активацией космической губки в парадигме IT3. Универсальное сигмоидальное масштабирование, наблюдаемое в данных ALICE, идентично перколяционному переходу сети войдов и филаментов.

Это позволяет предположить, что топология IT3 является не просто граничным условием для космологии, а фундаментальным геометрическим принципом, диктующим перенос энергии на всех масштабах физики. Вселенная — это единая, связанная, перколирующая система.

\section*{Благодарности}

Автор благодарит коллаборацию ALICE за открытый доступ к данным, который позволил провести это кросс-масштабное сравнение. Также выражается признательность теоретическим основам теории перколяции в статистической механике.

\begin{thebibliography}{99}

\bibitem[ALICE Collaboration(2026)]{ALICE2026}
ALICE Collaboration, 2026, Phys. Rev. Lett., 132, 042301

\bibitem[Logvinovich(2026a)]{Logvinovich2026a}
Logvinovich V., 2026a, Zenodo, DOI:10.5281/zenodo.19431323 (Статья I)

\bibitem[Logvinovich(2026c)]{Logvinovich2026c}
Logvinovich V., 2026c, Zenodo, DOI:10.5281/zenodo.19479551 (Статья III)

\bibitem[Logvinovich(2026d)]{Logvinovich2026d}
Logvinovich V., 2026d, Zenodo, DOI:10.5281/zenodo.19489307 (Статья IV)

\end{thebibliography}

\end{document}